%!TEX root = ../main.tex

% ============================================================
%  CAPÍTULO 3: MARCO CONCEPTUAL
% ============================================================

\chapter{Marco conceptual}\label{ch:marco}

% ──────────────────────────────────────────────────────────────
\section{\textit{Nowcasting} electoral y predicción con redes sociales}
% ──────────────────────────────────────────────────────────────

\subsection{Definición de \textit{nowcasting} en contexto electoral}

En ciencia política, \textit{nowcasting} (pronóstico del presente) designa una estimación cuantitativa de ``lo que pasaría si la elección se celebrara hoy''. Se diferencia del \textit{forecasting} clásico (predicción a futuro) en que el \textit{nowcasting} busca modelar el estado actual y latente del electorado mediante actualizaciones continuas y en tiempo real, basándose en un modelo cuyos parámetros estructurales se mantienen fijos mientras varían las covariables observadas \cite{lewisbeck2011nowcasting, lewisbeck2014proxy}.

La primera formulación formal aplicada a elecciones nacionales se consolidó con los modelos \textit{proxy}. Por ejemplo, en el Reino Unido, \citeauthor{lewisbeck2011nowcasting} (\citeyear{lewisbeck2011nowcasting}) modelaron el voto del oficialismo a partir de indicadores adelantados, como la aprobación del gobierno, medida con meses de antelación. En épocas más recientes, con el auge del \textit{Big Data}, la literatura empírica se ha volcado al uso de redes sociales como fuente masiva de covariables en tiempo real, transformando el \textit{nowcasting} en una herramienta de minería de datos dinámicos \cite{skoric2020}. 

Desde el estudio seminal de \citeauthor{tumasjan2010} (\citeyear{tumasjan2010}), múltiples autores han explorado arquitecturas híbridas. Trabajos recientes como los de \citeauthor{brito2022measuring} (\citeyear{brito2022measuring, brito2023somen}) han sofisticado la técnica combinando recolección de datos durante la campaña (\textit{During Campaign}) y modelos de aprendizaje automático (\textit{Machine Learning}) aplicados a Latinoamérica. Sin embargo, estas aplicaciones mantienen una dependencia crítica: el entrenamiento constante frente a encuestas tradicionales.

% ──────────────────────────────────────────────────────────────
\section{Modelos elección-agnósticos: definición y ventajas}
% ──────────────────────────────────────────────────────────────

Un modelo \textbf{elección-agnóstico} es aquel que se entrena utilizando datos conjuntos de múltiples elecciones simultáneamente, con el propósito de aprender reglas de decisión y pesos paramétricos que representen patrones de comportamiento electoral universales (o regionalmente consistentes), de modo que sean directamente transferibles a elecciones no vistas.

Formalmente, sea $\mathcal{E} = \{e_1, e_2, \ldots, e_M\}$ un conjunto de $M$ elecciones históricas diferentes (por ejemplo, alcaldías de distintas ciudades). El objetivo es aprender una función predictiva:

\begin{equation}
  f: \mathbb{R}^d \rightarrow [0,1]
  \label{eq:modelo-agnostico}
\end{equation}

\noindent tal que $\hat{y}_{i,m} = f(\mathbf{x}_{i,m})$, donde $\mathbf{x}_{i,m}$ es el vector de características de $d$ dimensiones (interacciones digitales, tasas de crecimiento, dominancia) del candidato $i$ en la elección $m$, y $\hat{y}_{i,m}$ es la cuota de voto estimada o la probabilidad de victoria.

La principal ventaja matemática y operativa de este enfoque frente a arquitecturas específicas de campaña como el SOMEN-DC \cite{brito2023somen} radica en su escalabilidad e independencia. Al abstraer la función $f$ sobre $M$ elecciones, el modelo se desacopla de la necesidad de contar con datos de encuestas como variable dependiente, utilizando directamente los resultados electorales finales como \textit{ground truth}. Asimismo, elimina la dependencia de infraestructura de \textit{scraping} en tiempo real, pues la predicción sobre una elección nueva $\mathcal{E}_{new}$ requiere únicamente extraer los vectores $\mathbf{x}_{new}$ y procesarlos a través del modelo pre-entrenado.

% ──────────────────────────────────────────────────────────────
\section{Interacciones digitales como \textit{proxy} de intención de voto}
% ──────────────────────────────────────────────────────────────

\subsection{La hipótesis Tumasjan}

El uso de señales digitales como estimador político se sustenta empíricamente en el trabajo fundacional de \citeauthor{tumasjan2010} (\citeyear{tumasjan2010}) sobre las elecciones federales alemanas de 2009. Estos autores demostraron que el simple conteo volumétrico de \textit{tweets} que mencionan a un partido refleja fielmente el resultado electoral, logrando un error absoluto medio (MAE) de 1.65\%. Esta hipótesis sugiere que la atención (reflejada en interacciones y contenido generado) equivale a capital político en la arena digital.

\subsection{Interacciones absolutas vs. dominancia relativa}

Si bien el volumen absoluto de interacciones (por ejemplo, el total de \textit{likes}) ofrece una señal importante, presenta ruido debido a la heterogeneidad de los distritos electorales y la presencia de cuentas atípicamente virales. Para mitigar esto, este estudio introduce y prioriza la métrica de \textbf{dominancia relativa}: el porcentaje de interacciones que acumula un candidato sobre el total de interacciones generadas por todos los candidatos competitivos dentro de su misma ciudad o distrito.

Los análisis exploratorios sobre las elecciones locales colombianas de 2023 presentados en este trabajo confirman que la dominancia relativa es significativamente más informativa que el volumen absoluto.

No obstante, se deben reconocer las limitaciones intrínsecas de este \textit{proxy}. Las redes sociales son ecosistemas permeables a la actividad inorgánica (bots, manipulación coordinada), sesgos demográficos propios de las plataformas y la existencia de candidatos que, por su condición de figuras públicas o \textit{influencers}, generan tracción digital desproporcionada a su verdadera viabilidad electoral.

% ──────────────────────────────────────────────────────────────
\section{El problema del histórico de interacciones}
% ──────────────────────────────────────────────────────────────

Uno de los mayores desafíos para entrenar modelos predictivos retrospectivos es la opacidad temporal de los datos en redes sociales. Plataformas como Facebook, Twitter/X y TikTok muestran únicamente el total acumulado de interacciones de una publicación en el instante en que es consultada (hoy). No guardan ni exponen públicamente la curva histórica de crecimiento; por lo tanto, es imposible observar directamente cuántos \textit{likes} tenía una publicación específicamente el día anterior a las elecciones, un dato crucial para el \textit{nowcasting}.

Para sortear este obstáculo metodológico, este trabajo propone modelar matemáticamente la curva de crecimiento acumulado de las interacciones utilizando una distribución de Weibull con parámetro de desplazamiento (\textit{offset}):

\begin{equation}
  F(t) = 1 - e^{-k \cdot (t + t_0)^{\alpha}}
  \label{eq:weibull}
\end{equation}

\noindent donde $t$ es el tiempo transcurrido (en días) desde la publicación; $k$ es el parámetro de escala que determina la velocidad general de acumulación; $\alpha$ es el parámetro de forma que captura la rápida deceleración típica del consumo de contenido en redes (decaimiento exponencial); y $t_0$ modela la asimetría temporal inicial, capturando la explosión masiva de interacciones que ocurre en las primeras horas antes de que $t$ se convierta en 1 día. 

Al calibrar estos parámetros ($k, \alpha, t_0$) de forma diferenciada para cada plataforma, es posible proyectar la curva $F(t)$ hacia el pasado y estimar con alta fidelidad las métricas que tenía cualquier publicación histórica en la víspera electoral.
